% Native Eq/Q/D bootstrap for the new B+monus arithmetic basis.
% Requires the core native rules to have already been explicitly stated.
\section{Native cases before a numerical equality characteristic}
\label{sec:monus-native-bootstrap}

There is a potential circularity in importing the preceding arithmetic
proof directly into the native Eq/$Q$/$D$ calculus. Its theorem-level
numerical case rule is not automatically a contextual split on $t=0$
and $t=Sz$ when the retained induction hypothesis mentions $t$. Moreover,
the tied-seed arithmetic has no unrestricted numerical selector among its
term functions. The following bounded selector removes that circularity.
Only hypothesis use, full formula transport, the native $Q/D$ rules,
the $B$ computation schemas, and literal-equation progress are used.
No ordinary implication, conjunction, disjunction, or quantifier rule is
assumed in this section.

\subsection{Initial native resources}
The native rules derive numerical cases and literal-equation induction
with their usual freshness conditions, using an unused witness in $D$,
then $D$ elimination and $Q$ specialization. The computation
$C(S0,p)=p$, used as both premises of transport with motive $w=p$,
gives $p=p$. Symmetry, transitivity, and numerical congruence follow
from the usual transport motives. The $B$ macros $P,Z,G$ of the previous
section are therefore available by their same equation derivations.

Equation explosion is also available. Under $0=1$, transport between
$C(0,a)=Sa$ and $C(1,a)=a$ gives $Sa=a$ for arbitrary $a$.
Literal-equation induction on $a=0$ then proves every $a=0$ in this
retained context. Hence every numerical equation follows. External
induction on a fixed native formula, using case introduction for each
$D$ and the unused-witness representation of $Q$, gives explosion for
all native formulas. This is an external finite proof transformation;
it is not compound-formula induction in the object theory.

For a numerical term $b$, say $\mathrm{FR}_{\{0,1\}}(b)$ holds when,
for every fixed native formula template $K$, the following metarule is
available in an arbitrary retained context:
\[
 \Gamma\vdash K(0),\quad\Gamma\vdash K(1)
       \quad\Longrightarrow\quad\Gamma\vdash K(b).
 \tag{FR}\label{monus:native-FR}
\]
The context may mention $b$; the premises add no equations $b=0$ or
$b=1$. Native numerical cases on the fresh template $w\mapsto K(Zw)$,
using $Z0=1$ and $Z(Sw)=0$, prove
$\mathrm{FR}_{\{0,1\}}(Zt)$ for every $t$. In particular,
$Z(Zb)=b$ follows for every such finite-range $b$ by applying FR to
that equation's two numeral instances.

\subsection{A selector whose first output is Boolean}
Define
\[
 J(a,b)=C(Zb,a),\qquad
 H(y,b,u)=J\bigl(G(y,u),G(Zy,b)\bigr).
 \tag{H}\label{monus:H}
\]
Direct computations give
\[
 J(a,0)=a,\quad J(a,1)=Sa,\qquad
 H(0,b,u)=b\quad(b\in\{0,1\}),\quad H(Sy,b,u)=u.
 \tag{H endpoints}\label{monus:H-endpoints}
\]
The first endpoint also holds for any term $b$ with the displayed
functional range rule: apply that rule to the fixed equation template.
The successor endpoint holds without any restriction on $b$.

This is not an unrestricted numerical selector. Its zero output is
Boolean; its successor output is arbitrary. Consequently its existence
does not contradict the one-coordinate majorant bound for tied-seed
terms. It is exactly enough for contextual Boolean case elimination.

\begin{lemma}[Contextual Boolean cases for equation conclusions]
\label{monus:contextual-Boolean}
Suppose $\mathrm{FR}_{\{0,1\}}(b)$ holds. If
\[
 \Gamma,b=0\vdash U=V,\qquad
 \Gamma,b=1\vdash U=V,
\]
then $\Gamma\vdash U=V$ is derivable in the native calculus. The retained
context and the conclusion may mention $b$.
\end{lemma}
\begin{proof}
With fresh $a,y$, define
\[
 K(x)=D_{a,y}\bigl(H(y,x,U)=H(y,b,V)\bigr),
\]
where $a$ is unused. For $r=0,1$, use case introduction to prove $K(r)$.
Its local zero assumption reduces to $r=b$, by the Boolean endpoint
of $H$ for $r$ and for $b$. Symmetry gives $b=r$, so cut in the
corresponding supplied branch proof. The successor matrix is $U=V$
by the other endpoint of $H$. This constructs $\Gamma\vdash K(r)$
for both numeral values, discharging the local equation.

Apply FR to obtain $\Gamma\vdash K(b)$. Paired elimination opens its
conditional closure. Its zero instance is $b=b$, supplied by
reflexivity and the endpoint computations. Specialization at $1$
returns $U=V$. All temporary variables are fresh, and all temporary
assumptions have been discharged.
\end{proof}

\subsection{Contextual zero/successor splits for arbitrary terms}
Native numerical cases, applied to fresh variables in the following
literal-equation templates, prove
\[
 G(Zt,t)=0,
 \qquad
 G(Z(Zt),t)=G(Z(Zt),S(Pt)).
 \tag{Reconstruction}\label{monus:native-reconstruction}
\]
For the first identity the zero branch is $G(1,0)=0$ and every
successor branch has control zero. For the second, the zero branch
has control zero, while the successor branch has control one and
$P(Sx)=x$. There is no substitution into a retained context.

\begin{lemma}[Contextual numerical cases for equations]
\label{monus:contextual-numerical}
Let $z$ be fresh for $\Gamma,t,U,V$. If
\[
 \Gamma,t=0\vdash U=V,\qquad
 \Gamma,t=Sz\vdash U=V,
\]
then $\Gamma\vdash U=V$ is derivable natively.
\end{lemma}
\begin{proof}
Apply Lemma~\ref{monus:contextual-Boolean} to $b=Zt$.
Under $Zt=1$, the first reconstruction identity yields $t=0$,
so the first given branch applies. Under $Zt=0$, the second identity
has control $Z0=1$, so it yields $t=S(Pt)$. Substitute $Pt$ for the
fresh $z$ in the second given branch and cut in this equation.
The conclusion is the same equation $U=V$ in both Boolean branches.
\end{proof}

Every case split in the gated-complement proof has an equation as its
ultimate conclusion. The zero cases that contradict $0\ne1$ are
replaced in the native calculus by the equation-explosion transformation
above. Thus the proof of the gated-complement invariant, antisymmetry,
and the two directions of equality reflection translates natively
without importing ordinary logical rules. Its only induction matrix
remains the single equation in the displayed invariant.

\subsection{Reconstructing ordinary logic after reflection}
The equality term $\epsilon(s,t)$ now has the native functional Boolean
range rule and satisfies
\[
 s=t\ \dashv\vdash\ \epsilon(s,t)=1
\]
in every retained context. Its finite range follows from the two
functional range rules for its $Z$ inputs and the gate endpoints.
For an atomic formula $E\equiv(s=t)$ define a gate by
\[
 \mathcal G_v(E)\;:\equiv\;
        G(v,\epsilon(s,t))=G(v,1).
\]
It is provable at $v=0$ and interderivable with $E$ at each $v=Sw$.
For atomic $E\equiv(s=t)$ and $F\equiv(r=w)$ define a formula selector
\[
 \mathcal M_v(E,F)\;:\equiv\;
        H(v,\epsilon(s,t),\epsilon(r,w))=1.
\]
The Boolean endpoint of $H$ gives $\mathcal M_0(E,F)\dashv\vdash E$;
its successor endpoint gives $\mathcal M_{Sz}(E,F)\dashv\vdash F$.
The zero branch is bounded, so no unrestricted numerical selector has
been introduced.

Extend the gate through $D$, replace $Q$ by its unused-witness $D$
representation, and align binder depths using the usual gate padding.
Then extend the formula selector through matched $D$ binders by external
recursion on their finite syntax. Binder congruence uses only the native
case rule, paired elimination, transport, and ordinary-assumption cut,
so it is already available. These are the same structural logical
constructions as in the inherited Choice-DI bridge~\cite{certified}, with
the atomic selector just given replacing its unbounded numerical selector.
Both inherited derivations are included under
\file{reference/initialized_original/} as
\file{inherited_choice_basis.tex} and
\file{inherited_ordinary_bridge.tex}.

For explicitness, put $p(z)=L(z;u,v.Zv)$, so $p(1)=1$ and
$p(2)=0$. With all unused coordinates fresh, the ordinary macros are
\begin{align*}
 A\Rightarrow_*B&=D_{a,z}\mathcal M_z(A,B),\\
 A\wedge_*B&=D_{a,z}\mathcal G_z\mathcal M_{p(z)}(A,B),\\
 \exists_*x\,A(x)&=D_{x,y}\mathcal G_y A(x),\\
 A\vee_*B&=\exists_*x\,\mathcal M_x(A,B),\\
 \forall_*x\,A(x)&=A(0)\wedge_*D_{a,z}\mathcal G_z A(z),\\
 \bot_*&=(0=1).
\end{align*}
Implication introduction uses the matrix's zero assumption; elimination
opens its closure and specializes at one. The gate's provable zero
instance supplies conjunction and existential introduction, while
specialization at one or two supplies their eliminations. After implication
is available, ordinary contextual numerical cases for arbitrary formulas
follow by applying Native NC to an implication template; these cases
justify disjunction elimination and universal elimination. All freshness
and capture-avoidance conditions are the inherited ones.

These standard native macros for implication, conjunction, disjunction,
existence, universal quantification, and falsity therefore have their usual
intuitionistic introduction and elimination rules. Both formula
composites follow from the same gate and selector endpoint proofs.
After this stage the ordinary $T_{Bd}$ argument and the inherited
bounded compiler may be used without a logical bootstrap cycle.
